\documentclass[11pt]{article}
\usepackage[T1]{fontenc}
\usepackage[utf8]{inputenc}
\usepackage{lmodern}
\usepackage[margin=1in]{geometry}
\usepackage{amsmath,amssymb,amsthm,mathtools}
\usepackage{enumitem}
\usepackage[hidelinks]{hyperref}

\newtheorem{theorem}{Theorem}[section]
\newtheorem{proposition}[theorem]{Proposition}
\newtheorem{definition}[theorem]{Definition}
\newtheorem{lemma}[theorem]{Lemma}
\newtheorem{corollary}[theorem]{Corollary}
\newtheorem{remark}[theorem]{Remark}
\newtheorem{example}[theorem]{Example}

\title{Star and Support-Image Profiles for Class-Two $p$-Groups}
\author{Luca Blanchi}
\date{June 2026}

\begin{document}
\maketitle

\begin{abstract}
Let $G$ be a finite group. The higher commuting probabilities
\[
P_m(G)=\Pr(g_1,\dots,g_m\text{ commute pairwise})=\frac{|\operatorname{Hom}(\mathbb Z^m,G)|}{|G|^m}
\]
measure the abundance of abelian $m$-tuples in $G$. We compare these probabilities with several other natural quadratic observers. The first family is the star commuting hierarchy
\[
\Sigma_{r,t}(G)=\Pr([x_i,y_j]=1\ \forall i,j)=\frac{|\operatorname{Hom}(F_r\times F_t,G)|}{|G|^{r+t}}.
\]
We show that $\Sigma_{r,t}(G)$ is a moment of common-centralizer sizes. For groups of class $2$ and exponent $p$, represented by alternating maps $\beta:V\times V\to W$, the star hierarchy becomes linear algebra: if $A_{\mathbf v}:V\to W^r$ is the common contraction map and $N_r^\beta(a)$ counts $r$-tuples for which $\operatorname{rank}A_{\mathbf v}=a$, then $\Sigma_{r,t}(G_\beta)=|V|^{-r}\sum_aN_r^\beta(a)p^{-ta}$. By contrast, $P_m(G_\beta)$ is determined by the numbers of totally isotropic subspaces of $V$. We prove that these two kinds of information are genuinely distinct by constructing, for every odd prime $p$, two directly indecomposable special groups $G_A(p)$ and $G_B(p)$ of order $p^7$, class $2$, and exponent $p$, with identical $P_m$ for all $m$ but different $N_1$. We then identify the star hierarchy with rank-support enumerators of the rank-metric code of contractions, prove a positive completeness theorem for complete multipartite graph blow-up groups, introduce canonical support-image profiles $\operatorname{CSI}_r$, and close with the Centralizer Hom Transform.
\end{abstract}

\section*{Declaration of generative AI and AI-assisted technologies in the writing process}
This article belongs to the AI-assisted math section. Generative AI, including ChatGPT by \mbox{OpenAI}, was used to assist with mathematical drafting, formalization, review, and editing. The note may be preliminary, may have been only partially reviewed, and in some cases may be only AI-reviewed.

\section{Introduction}

The commuting probability of a finite group is

\[
\operatorname{cp}(G)
=
\Pr([x,y]=1)
=
\frac{|\operatorname{Hom}(\mathbb Z^2,G)|}{|G|^2}.
\]

A natural hierarchy of refinements is given by

\[
P_m(G)
=
\Pr(g_1,\dots,g_m\text{ commute pairwise})
=
\frac{|\operatorname{Hom}(\mathbb Z^m,G)|}{|G|^m}.
\]

These are the higher commuting probabilities. They count homomorphisms from free abelian groups.

This paper studies a different refinement. Instead of asking for all elements in a tuple to commute pairwise, we ask for two tuples to commute across:

\[
[x_i,y_j]=1
\qquad
(1\le i\le r,\ 1\le j\le t).
\]

The resulting probabilities are

\[
\Sigma_{r,t}(G)
=
\Pr([x_i,y_j]=1\ \forall i,j)
=
\frac{|\operatorname{Hom}(F_r\times F_t,G)|}{|G|^{r+t}}.
\]

Thus $P_m$ counts homomorphisms from $\mathbb Z^m$, while $\Sigma_{r,t}$ counts homomorphisms from direct products of free groups.

The main contribution of the paper is not a total ordering of all these invariants. Rather, we compare several natural families of quadratic observers and show that they retain different aspects of the alternating geometry of a class-two exponent-$p$ group.

Such groups are governed by alternating bilinear maps

\[
\beta:V\times V\to W.
\]

For $G_\beta$, the higher commuting probabilities $P_m$ are controlled by totally isotropic subspaces of $V$. The star probabilities $\Sigma_{r,t}$, on the other hand, are controlled by the ranks of common contraction maps

\[
A_{\mathbf v}:V\to W^r.
\]

These two data sets are genuinely different. The uniform family constructed in Section 6 shows that $P_m$ for all $m$ does not determine $N_1$, even on directly indecomposable groups.

The rank side has an independent algebraic meaning. The contraction maps

\[
A_v:u\mapsto \beta(v,u)
\]

form a rank-metric code

\[
C_\beta\le \operatorname{Hom}(V,W),
\]

and the profiles $N_r$ are equivalent to the higher rank-support enumerators of this code. Thus star probabilities connect finite-group word observers to rank-metric support theory and $q$-polymatroidal invariants.

The canonical support-image profile $\operatorname{CSI}_r$ is introduced as a further codomain-localized refinement. It remembers the $\operatorname{GL}(W)$-orbit of the image subspace

\[
\operatorname{im}A_{\mathbf v}\le W^r,
\]

rather than only its dimension or the dimensions of its coordinate projections.

The Centralizer Hom Transform is different in nature. It does not refine CSI in any formal sense asserted here; instead, it replaces centralizer-size moments by homomorphism counts into centralizers, and therefore records internal isomorphism types of common centralizers.


\section{Star commuting probabilities}

Let $G$ be a finite group. For $r,t\ge0$, define

\[
\Sigma_{r,t}(G)=
\Pr([x_i,y_j]=1\ \forall i,j),
\]

where

\[
x_1,\dots,x_r,y_1,\dots,y_t
\]

are independent uniformly distributed elements of $G$. We use the convention

\[
\Sigma_{r,0}(G)=\Sigma_{0,t}(G)=1.
\]

Let

\[
\mathbf x=(x_1,\dots,x_r)\in G^r.
\]

Write

\[
C_G(\mathbf x)
=\bigcap_{i=1}^rC_G(x_i).
\]

\begin{proposition}[Centralizer moment formula]

For every finite group $G$,

\[
\Sigma_{r,t}(G)
=\frac1{|G|^r}
\sum_{\mathbf x\in G^r}
\left(
\frac{|C_G(\mathbf x)|}{|G|}
\right)^t.
\]

\end{proposition}
\begin{proof}

Fix $\mathbf x=(x_1,\dots,x_r)$. The condition

\[
[x_i,y_j]=1
\qquad
\forall i,j
\]

is equivalent to

\[
y_j\in C_G(\mathbf x)
\qquad
\forall j.
\]

Thus, for fixed $\mathbf x$, the conditional probability is

\[
\left(
\frac{|C_G(\mathbf x)|}{|G|}
\right)^t.
\]

Averaging over all $\mathbf x\in G^r$ gives the formula.
\end{proof}

A homomorphism

\[
F_r\times F_t\to G
\]

is precisely a pair of tuples whose entries commute across the two factors. Hence

\[
\Sigma_{r,t}(G)
=\frac{|\operatorname{Hom}(F_r\times F_t,G)|}{|G|^{r+t}}.
\]

The case $r=t=1$ is the usual commuting probability.

\section{Class-two exponent-(p) groups}

Let $p$ be an odd prime. Let $V,W$ be finite-dimensional vector spaces over $\mathbb F_p$, and let

\[
\beta:V\times V\to W
\]

be an alternating bilinear map.

Define

\[
G_\beta=V\oplus W
\]

with multiplication

\[
(v,z)(u,w)
\left(
v+u,
z+w+\frac12\beta(v,u)
\right).
\]

Then

\[
[(v,z),(u,w)]
(0,\beta(v,u)).
\]

Thus $G_\beta$ has nilpotency class at most $2$ and exponent $p$.

For

\[
v\in V,
\]

define the contraction

\[
A_v:V\to W,
\qquad
A_v(u)=\beta(v,u).
\]

For an $r$-tuple

\[
\mathbf v=(v_1,\dots,v_r)\in V^r,
\]

define

\[
A_{\mathbf v}:V\to W^r,
\qquad
u\mapsto
(A_{v_1}u,\dots,A_{v_r}u).
\]

Let

\[
N_r^\beta(a)=
|{\mathbf v\in V^r:\operatorname{rank}A_{\mathbf v}=a}|.
\]

\begin{proposition}[Star probabilities as rank moments]

For $G_\beta$,

\[
\Sigma_{r,t}(G_\beta)=
|V|^{-r}
\sum_a
N_r^\beta(a)p^{-ta}.
\]

\end{proposition}
\begin{proof}

Let

\[
x_i=(v_i,z_i)\in G_\beta.
\]

The common centralizer of the tuple $x_1,\dots,x_r$ consists of all

\[
(u,w)\in V\oplus W
\]

such that

\[
\beta(v_i,u)=0
\qquad
\forall i.
\]

Thus

\[
|C_{G_\beta}(x_1,\dots,x_r)|
=|W|\cdot |\ker A_{\mathbf v}|
=|W|\cdot p^{\dim V-\operatorname{rank}A_{\mathbf v}}.
\]

Since

\[
|G_\beta|=|V||W|,
\]

we obtain

\[
\frac{|C_{G_\beta}(x_1,\dots,x_r)|}{|G_\beta|}
p^{-\operatorname{rank}A_{\mathbf v}}.
\]

The central coordinates $z_i$ do not affect the rank. Substituting into the centralizer moment formula gives

\[
\Sigma_{r,t}(G_\beta)=
|V|^{-r}
\sum_{\mathbf v\in V^r}
p^{-t\operatorname{rank}A_{\mathbf v}},
\]

which is the stated formula.
\end{proof}

Thus the values $\Sigma_{r,t}$, as $t$ varies, are precisely the moments of the rank distribution $N_r^\beta$.

\section{Higher commuting probabilities and isotropic subspaces}

A subspace

\[
U\le V
\]

is totally isotropic for $\beta$ if

\[
\beta(U,U)=0.
\]

Let

\[
I_d(\beta)
\]

be the number of $d$-dimensional totally isotropic subspaces of $V$.

\begin{proposition}[$P_m$ and isotropic subspaces]

For $G_\beta$,

\[
P_m(G_\beta)
=|V|^{-m}
\sum_d
I_d(\beta)
\prod_{i=0}^{d-1}(p^m-p^i).
\]

\end{proposition}
\begin{proof}

An $m$-tuple in $G_\beta$ commutes pairwise if and only if its image in $V$ spans a totally isotropic subspace.

Fix a $d$-dimensional totally isotropic subspace $U\le V$. The number of ordered $m$-tuples spanning $U$ is the number of surjective linear maps

\[
\mathbb F_p^m\twoheadrightarrow U,
\]

namely

\[
\prod_{i=0}^{d-1}(p^m-p^i).
\]

The central coordinates in $W$ are arbitrary and cancel in the normalization by $|G_\beta|^m$. Summing over all totally isotropic $U$ gives the formula.
\end{proof}

Therefore the entire hierarchy $P_m$ depends only on the isotropic-subspace counts

\[
I_d(\beta).
\]

The star hierarchy sees different data: ranks of common contraction maps.

\section{Rank-support enumerators}

Assume in this section that the radical of $\beta$ is zero, so that the map

\[
V\to\operatorname{Hom}(V,W),
\qquad
v\mapsto A_v
\]

is injective. Let

\[
C_\beta={A_v:v\in V}\le\operatorname{Hom}(V,W)
\]

be the corresponding rank-metric code of contractions.

For a subcode

\[
D\le C_\beta,
\]

define its row-support weight by

\[
\operatorname{wt}_R(D)=
\dim V-\dim\bigcap_{A\in D}\ker A.
\]

For $j,a\ge0$, define

\[
A_{j,a}(C_\beta)=
|{D\le C_\beta:\dim D=j,\ \operatorname{wt}_R(D)=a}|.
\]

\begin{theorem}[Star profiles and rank-support enumerators]

For all $r,a$,

\[
N_r^\beta(a)=
\sum_{j=0}^{\min(r,\dim C_\beta)}
\gamma_{r,j}(p)A_{j,a}(C_\beta),
\]

where

\[
\gamma_{r,j}(p)
=\prod_{i=0}^{j-1}(p^r-p^i).
\]

Consequently, the family

\[
(N_r^\beta(a))_{r,a}
\]

determines, and is determined by, the family

\[
(A_{j,a}(C_\beta))_{j,a}.
\]

\end{theorem}
\begin{proof}

For an $r$-tuple

\[
\mathbf v=(v_1,\dots,v_r),
\]

let

\[
D(\mathbf v)
=\langle A_{v_1},\dots,A_{v_r}\rangle
\le C_\beta.
\]

Then

\[
\operatorname{rank}A_{\mathbf v}
=\dim V-\dim\bigcap_{i=1}^r\ker A_{v_i}
=\operatorname{wt}_R(D(\mathbf v)).
\]

Fix a $j$-dimensional subcode $D\le C_\beta$. The number of ordered $r$-tuples generating $D$ is the number of surjective linear maps

\[
\mathbb F_p^r\twoheadrightarrow D,
\]

namely

\[
\gamma_{r,j}(p)=
\prod_{i=0}^{j-1}(p^r-p^i).
\]

Summing over all $D$ of dimension $j$ and row-support weight $a$ gives the stated formula.

The matrix ($\gamma_{r,j}$) is triangular in $r,j$, with nonzero diagonal

\[
\gamma_{j,j}(p)=|\operatorname{GL}_j(\mathbb F_p)|.
\]

Hence the relation is invertible.
\end{proof}

This identifies the star hierarchy with higher rank-support enumerators of the contraction code.

\section{An infinite separation between $P_m$ and $N_1$}

We now show that the higher commuting hierarchy $P_m$ does not determine even the first star profile $N_1$.

Let

\[
V=\mathbb F_p^4
\]

with basis

\[
e_1,e_2,e_3,e_4.
\]

Use the ordered basis

\[
e_{12},e_{13},e_{14},e_{23},e_{24},e_{34}
\]

of $\Lambda^2V$.

Define subspaces

\[
K_A=
\langle
e_{12},\ e_{34},\ e_{13}+e_{24}
\rangle
\]

and

\[
K_B=
\langle
e_{12},\ e_{14},\ e_{13}+e_{24}
\rangle.
\]

Let

\[
W_A=\Lambda^2V/K_A,
\qquad
W_B=\Lambda^2V/K_B,
\]

and let

\[
\beta_A,\beta_B:\Lambda^2V\to W_A,W_B
\]

be the quotient maps.

\begin{lemma}[The $K_A$-isotropic planes are disjoint outside zero]

Let

\[
K_A=\langle e_{12},e_{34},e_{13}+e_{24}\rangle.
\]

The $p+1$ two-dimensional subspaces of $V$ that are totally isotropic for $\beta_A$ intersect pairwise only in $0$.

\end{lemma}
\begin{proof}

A two-dimensional subspace $U\le V$ is totally isotropic for $\beta_A$ precisely when its Plucker point lies in

\[
\mathbb P(K_A)\cap\mathcal K,
\]

where $\mathcal K$ is the Klein quadric.

For $K_A$, this intersection is the nondegenerate conic

\[
ac-b^2=0
\]

in $\mathbb P(K_A)\cong\mathbb P^2$. In particular, it contains no projective line.

Suppose two distinct isotropic planes $U_1,U_2\le V$ intersect in a nonzero vector $v$. Choose vectors $u_1,u_2$ such that

\[
U_i=\langle v,u_i\rangle.
\]

Then the Plucker points

\[
[v\wedge u_1],\qquad [v\wedge u_2]
\]

are two distinct points of

\[
\mathbb P(K_A)\cap\mathcal K.
\]

Moreover, the projective line joining them is

\[
\{[v\wedge(\lambda u_1+\mu u_2)]:[\lambda:\mu]\in\mathbb P^1\}.
\]

Every point on this line is decomposable, hence lies on the Klein quadric. Since the endpoints lie in $\mathbb P(K_A)$, the whole line lies in $\mathbb P(K_A)\cap\mathcal K$.

This contradicts the fact that $\mathbb P(K_A)\cap\mathcal K$ is a nondegenerate conic and contains no line. Therefore distinct isotropic planes intersect only in $0$.
\end{proof}

\begin{theorem}[Infinite directly indecomposable separation]

For every odd prime $p$,

\[
P_m(G_{\beta_A})=P_m(G_{\beta_B})
\qquad
\forall m\ge1,
\]

but

\[
N_1^{\beta_A}\ne N_1^{\beta_B}.
\]

Moreover, both groups $G_{\beta_A}$ and $G_{\beta_B}$ are directly indecomposable.

\end{theorem}
\begin{proof}

The planes in $V$ are parametrized by the Klein quadric in

\[
\mathbb P(\Lambda^2V),
\]

given in Plucker coordinates by

\[
x_{12}x_{34}-x_{13}x_{24}+x_{14}x_{23}=0.
\]

A plane is totally isotropic for $\beta_X$ if and only if its Plucker point lies in

\[
\mathbb P(K_X)
\]

and on the Klein quadric.

For $K_A$, a vector has coordinates

\[
(a,b,0,0,b,c),
\]

and the Klein equation restricts to

\[
ac-b^2=0.
\]

This is a nondegenerate conic in $\mathbb P^2(\mathbb F_p)$, hence has $p+1$ points.

For $K_B$, a vector has coordinates

\[
(a,b,c,0,b,0),
\]

and the Klein equation restricts to

\[
-b^2=0.
\]

This is a line in $\mathbb P^2(\mathbb F_p)$, again with $p+1$ points.

Thus $\beta_A$ and $\beta_B$ have the same number of totally isotropic planes. Every one-dimensional subspace is isotropic because $\beta$ is alternating. Since there are only $p+1$ isotropic planes, no three-dimensional subspace can be totally isotropic, as any three-dimensional isotropic subspace would contain more than $p+1$ planes.

Therefore the full sequence

\[
I_d(\beta_A)
\]

agrees with

\[
I_d(\beta_B).
\]

By Proposition 4.1,

\[
P_m(G_{\beta_A})=P_m(G_{\beta_B})
\]

for all $m$.

It remains to compute $N_1$. For $0\ne v\in V$,

\[
\operatorname{rank}A_v
=3-\dim(K_X\cap(v\wedge V)),
\]

because $v\wedge V$ has dimension $3$ and $W_X=\Lambda^2V/K_X$.

By Lemma 6.2, the nonzero vectors lying in the isotropic planes for $K_A$ are exactly $(p+1)(p^2-1)$. Equivalently, the intersection of $\mathbb P(K_A)$ with the Klein quadric is a nondegenerate conic, which contains no projective line. Hence no nonzero $v$ gives a two-dimensional intersection

\[
K_A\cap(v\wedge V).
\]

Thus there are no rank-one contractions.

The $p+1$ isotropic planes are pairwise disjoint outside $0$; otherwise the corresponding points of the conic would lie on a projective line contained in the Klein quadric and in $\mathbb P(K_A)$. Therefore the nonzero vectors lying in isotropic planes number

\[
(p+1)(p^2-1).
\]

For these vectors the contraction rank is $2$, and for the remaining nonzero vectors it is $3$. Hence

\[
(c_0,c_1,c_2,c_3)_A
\left(
1,
0,
(p+1)(p^2-1),
p(p-1)^2(p+1)
\right).
\]

For $K_B$, the $p+1$ isotropic planes form a pencil through the line

\[
\langle e_1\rangle.
\]

The $p-1$ nonzero vectors on this line have contraction rank $1$. The vectors lying in one of the isotropic planes but not on the common line have contraction rank $2$, and their number is

\[
(p+1)(p^2-p)=p(p^2-1).
\]

The remaining nonzero vectors have rank $3$. Thus

\[
(c_0,c_1,c_2,c_3)_B
\left(
1,
p-1,
p(p^2-1),
p^3(p-1)
\right).
\]

Therefore $N_1^{\beta_A}\ne N_1^{\beta_B}$.

Finally, we prove direct indecomposability. Both maps are surjective and radical-free, so the corresponding groups are special. A direct product decomposition of a special class-two group induces an orthogonal decomposition

\[
V=U_1\oplus U_2
\]

of the alternating map, with each factor contributing nontrivially to the commutator image. No $U_i$ can be one-dimensional. Since $\dim V=4$, any nontrivial such decomposition would have

\[
\dim U_1=\dim U_2=2.
\]

But then each $\beta(\Lambda^2U_i)$ has dimension at most $1$, so the total commutator image has dimension at most $2$, contradicting

\[
\dim W_A=\dim W_B=3.
\]

Thus both groups are directly indecomposable.
\end{proof}

This proves a uniform separation

\[
P_m<N_1
\]

on directly indecomposable groups.

\section{Complete multipartite blow-ups}

We now give a positive classification result.

Let

\[
V=V_1\oplus\cdots\oplus V_k,
\qquad
\dim V_i=n_i,
\qquad
k\ge2.
\]

Let

\[
W=\bigoplus_{1\le i<j\le k}V_i\otimes V_j.
\]

Define

\[
\beta(v,u)_{ij}
=v_i\otimes u_j-u_i\otimes v_j.
\]

This is the alternating map associated with the complete multipartite blow-up.

Let

\[
n=\sum_i n_i.
\]

\begin{lemma}[Rank formula]

For $v\in V$,

\[
\operatorname{rank}A_v=
\begin{cases}
0,&v=0,\\
n-n_i,&\operatorname{supp}(v)=\{i\},\\
n-1,&|\operatorname{supp}(v)|\ge2.
\end{cases}
\]

\end{lemma}
\begin{proof}

If $v=0$, the contraction is zero.

Suppose $v$ is supported only in $V_i$. Then

\[
A_v(u)=0
\]

if and only if

\[
u_j=0
\qquad
(j\ne i).
\]

Thus

\[
\ker A_v=V_i
\]

and

\[
\operatorname{rank}A_v=n-n_i.
\]

Now suppose $v$ has nonzero components in at least two blocks. If $j\notin\operatorname{supp}(v)$, then for some $i\in\operatorname{supp}(v)$, the edge component gives

\[
v_i\otimes u_j=0,
\]

hence

\[
u_j=0.
\]

If $i,j\in\operatorname{supp}(v)$, the equation

\[
v_i\otimes u_j=u_i\otimes v_j
\]

forces

\[
u_i=\lambda v_i,
\qquad
u_j=\lambda v_j
\]

with the same scalar $\lambda$. Therefore

\[
u=\lambda v.
\]

Hence

\[
\ker A_v=\langle v\rangle
\]

and

\[
\operatorname{rank}A_v=n-1.
\]

\end{proof}

\begin{theorem}[Completeness of $N_1$ for complete multipartite blow-ups]

For complete multipartite blow-ups, $N_1$ determines the multiset

\[
{n_1,\dots,n_k}.
\]

Consequently, $N_1$ determines the corresponding group.

\end{theorem}
\begin{proof}

From

\[
\sum_aN_1(a)=|V|=p^n
\]

we recover $n$.

For $s\ge2$, let

\[
m_s=|{i:n_i=s}|.
\]

By the rank formula, the number of nonzero vectors supported in a block of dimension $s$ is

\[
m_s(p^s-1),
\]

and these vectors have rank

\[
n-s.
\]

Thus

\[
m_s=
\frac{N_1(n-s)}{p^s-1}
\qquad
(s\ge2).
\]

Finally,

\[
m_1
=n-\sum_{s\ge2}s\,m_s.
\]

Hence the multiset of block dimensions is recovered.

The alternating map is determined by this multiset up to equivalence, so the group is determined.
\end{proof}

\begin{proposition}[Direct indecomposability]

Every complete multipartite blow-up group with $k\ge2$ is directly indecomposable.

\end{proposition}
\begin{proof}

A direct product decomposition would induce an orthogonal decomposition

\[
V=U_1\oplus U_2.
\]

If a nonzero vector of $U_1$ has support in at least two blocks, then by Lemma 7.1 its contraction kernel is one-dimensional. Orthogonality would force

\[
U_2\subseteq \ker A_v=\langle v\rangle,
\]

which is impossible in a nontrivial decomposition.

Thus every nonzero vector in $U_1$ is supported in a single block. By linearity, $U_1$ is contained in some $V_i$. Similarly, $U_2$ is contained in some $V_j$. If $i\ne j$, then

\[
\beta(U_1,U_2)\ne0.
\]

If $i=j$, then

\[
U_1+U_2
\]

does not contain the other blocks. Since $k\ge2$, this cannot be all of $V$. Contradiction.
\end{proof}

Thus $N_1$ is complete on a natural infinite directly indecomposable class.

\section{Flag profiles}

The profile $N_r$ records only the final rank of

\[
A_{\mathbf v}:V\to W^r.
\]

It forgets how ranks are distributed among subtuples.

For

\[
\mathbf v=(v_1,\dots,v_r)
\]

and nonempty

\[
S\subseteq [r],
\]

define

\[
\rho_{\mathbf v}(S)
=\operatorname{rank}A_{\mathbf v_S}.
\]

The aggregated flag profile of level $r$ is the distribution of the functions

\[
S\mapsto \rho_{\mathbf v}(S)
\]

as $\mathbf v$ ranges over $V^r$.

For $r=2$, this is the collection of numbers

\[
M_{a,b,c}
=|\{(v,w):
\operatorname{rank}A_v=a,\
\operatorname{rank}A_w=b,\
\operatorname{rank}A_{v,w}=c
\}|.
\]

Flag profiles refine $N_r$, but they remain aggregated: they remember only dimensions of projected images, not the actual position of the image subspaces in $W^r$.

The next section introduces a compact codomain-localized refinement.

\section{Canonical support-image profiles}

For

\[
\mathbf v=(v_1,\dots,v_r)\in V^r,
\]

define

\[
I_\beta(\mathbf v)=\operatorname{im}A_{\mathbf v}\le W^r.
\]

For a subspace

\[
T\le W^r,
\]

define

\[
E_r^\beta(T)
=|\{\mathbf v\in V^r:I_\beta(\mathbf v)=T\}|.
\]

The group $\operatorname{GL}(W)$ acts diagonally on $W^r$, hence on its subspaces.

\begin{definition}[CSI]

The canonical support-image profile of level $r$, denoted

\[
\operatorname{CSI}_r(\beta),
\]

is the function on orbits

\[
\mathcal O\in \operatorname{Sub}(W^r)/GL(W)
\]

given by

\[
\operatorname{CSI}_r(\beta)(\mathcal O)
=\sum_{T\in\mathcal O}E_r^\beta(T).
\]

For $R\ge1$, define

\[
\operatorname{CSI}_{\le R}(\beta)
(\operatorname{CSI}_1(\beta),\dots,\operatorname{CSI}_R(\beta)).
\]

\end{definition}
\begin{proposition}[Invariance]

If $\beta$ and $\beta'$ are equivalent under

\[
GL(V)\times GL(W),
\]

then

\[
\operatorname{CSI}_r(\beta)=\operatorname{CSI}_r(\beta')
\]

for all $r$.

\end{proposition}
\begin{proof}

Suppose

\[
\beta'(Pv,Pu)=Q\beta(v,u)
\]

with

\[
P\in GL(V),
\qquad
Q\in GL(W).
\]

Then

\[
I_{\beta'}(P\mathbf v)
=Q^{\oplus r}I_\beta(\mathbf v).
\]

The map

\[
\mathbf v\mapsto P\mathbf v
\]

is a bijection of $V^r$, and $Q^{\oplus r}$ is precisely the diagonal action of $\operatorname{GL}(W)$ on $W^r$. Thus the orbit distribution is preserved.
\end{proof}

\begin{proposition}[CSI dominates $N_r$ and flag profiles]

The profile $\operatorname{CSI}_r(\beta)$ determines $N_r^\beta$ and the aggregated flag profile of level $r$.

\end{proposition}
\begin{proof}

The rank profile $N_r$ is obtained by summing $\operatorname{CSI}_r$ over orbits of subspaces $T\le W^r$ of a fixed dimension:

\[
N_r^\beta(a)=
\sum_{\substack{\mathcal O\ \dim T=a}}
\operatorname{CSI}_r(\beta)(\mathcal O).
\]

For flag profiles, let

\[
\pi_S:W^r\to W^S
\]

be coordinate projection. For

\[
T=I_\beta(\mathbf v),
\]

we have

\[
\dim \pi_S(T)
=\rho_{\mathbf v}(S).
\]

The function

\[
T\mapsto
(\dim\pi_S(T))_{\varnothing\ne S\subseteq[r]}
\]

is invariant under diagonal $\operatorname{GL}(W)$. Hence the aggregated distribution of flag functions is determined by the CSI orbit distribution.
\end{proof}

Thus CSI is a codomain-localized refinement of both rank and flag profiles.

\section{Computing CSI}

Let

\[
d=\dim W.
\]

We now show that $\operatorname{CSI}_r$ is fixed-parameter computable in $d$ and $r$.

For

\[
S\le W^r,
\]

define

\[
F_r^\beta(S)
=|\{\mathbf v\in V^r:I_\beta(\mathbf v)\le S\}|.
\]

\begin{lemma}

For every $S\le W^r$, $F_r^\beta(S)$ is computable by linear algebra in time polynomial in $\dim V$, for fixed $d,r$.

\end{lemma}
\begin{proof}

Let

\[
\pi_S:W^r\to W^r/S
\]

be the quotient map. The condition

\[
I_\beta(\mathbf v)\le S
\]

is equivalent to

\[
\pi_S\circ A_{\mathbf v}=0.
\]

The map

\[
\Phi_S:V^r\to\operatorname{Hom}(V,W^r/S),
\qquad
\mathbf v\mapsto \pi_S\circ A_{\mathbf v}
\]

is linear. Hence

\[
F_r^\beta(S)=p^{\dim\ker\Phi_S},
\]

and the dimension of the kernel is computed by Gaussian elimination.
\end{proof}

\begin{theorem}[CSI computation]

For fixed $d=\dim W$ and $r$, $\operatorname{CSI}_r(\beta)$ is computable in time

\[
p^{O(d^2r^2)}\operatorname{poly}(\dim V,\log p).
\]

\end{theorem}
\begin{proof}

The number of subspaces of $W^r$, where

\[
\dim W^r=dr,
\]

is

\[
p^{O(d^2r^2)}.
\]

For each subspace $S\le W^r$, compute

\[
F_r^\beta(S)
\]

using Lemma 10.1.

We have

\[
F_r^\beta(S)
=\sum_{T\le S}E_r^\beta(T).
\]

By Mobius inversion on the lattice of subspaces,

\[
E_r^\beta(T)
=\sum_{S\le T}
\mu(S,T)F_r^\beta(S),
\]

where

\[
\mu(S,T)
=(-1)^c p^{\binom c2},
\qquad
c=\dim T-\dim S.
\]

This computes all values $E_r^\beta(T)$. Finally, aggregate them over diagonal $\operatorname{GL}(W)$-orbits on subspaces $T\le W^r$. Since the relevant ambient dimension is $dr$, this orbit aggregation also costs

\[
p^{O(d^2r^2)}.
\]

\end{proof}

CSI is therefore a compact substitute for full localization on the domain side. The full profile over all subspaces of $V$ has size roughly

\[
p^{\Theta((\dim V)^2)},
\]

whereas CSI$_r$ is controlled by the small space $W^r$.

\section{Centralizer Hom Transform}

We close with a further refinement, which sees not only sizes of centralizers but their isomorphism types.

Let $K$ be a finitely generated group. Define

\[
\operatorname{CHT}_r(G;K)
=|\operatorname{Hom}(F_r\times K,G)|.
\]

\begin{proposition}[CHT formula]

For every finite group $G$,

\[
\operatorname{CHT}_r(G;K)=
\sum_{\mathbf x\in G^r}
|\operatorname{Hom}(K,C_G(\mathbf x))|.
\]

\end{proposition}
\begin{proof}

A homomorphism

\[
F_r\times K\to G
\]

is determined by the images of the $r$ free generators of $F_r$, say

\[
\mathbf x=(x_1,\dots,x_r),
\]

together with a homomorphism

\[
K\to G
\]

whose image commutes with every $x_i$. Such a homomorphism has image in

\[
C_G(\mathbf x).
\]

Summing over $\mathbf x\in G^r$ gives the formula.
\end{proof}

For $K=F_t$, one recovers

\[
\operatorname{CHT}_r(G;F_t)
=\sum_{\mathbf x\in G^r}|C_G(\mathbf x)|^t
=|G|^{r+t}\Sigma_{r,t}(G).
\]

Thus the star hierarchy is the free-group part of CHT.

\begin{theorem}[CHT inversion]

Let $G$ be finite of order $n$. For fixed $r$, the values

\[
\operatorname{CHT}_r(G;K),
\qquad |K|\le n,
\]

with $K$ finite, determine the multiset

\[
{C_G(\mathbf x):\mathbf x\in G^r}/\cong
\]

with multiplicities.

\end{theorem}
\begin{proof}

Let

\[
\mu_{r,G}(H)
=|\{\mathbf x\in G^r:C_G(\mathbf x)\cong H\}|.
\]

Then

\[
\operatorname{CHT}_r(G;K)=
\sum_H
\mu_{r,G}(H)|\operatorname{Hom}(K,H)|.
\]

For finite groups $K,H$, write

\[
h_K(H)=|\operatorname{Hom}(K,H)|,
\]

and

\[
i_K(H)=|\operatorname{Inj}(K,H)|.
\]

Every homomorphism $K\to H$ has kernel $N\trianglelefteq K$ and induces an embedding

\[
K/N\hookrightarrow H.
\]

Therefore

\[
h_K(H)=
\sum_{N\trianglelefteq K}
i_{K/N}(H).
\]

By induction on quotients of $K$, the CHT values determine

\[
\sum_H \mu_{r,G}(H)i_K(H)
\]

for every $K$.

Now order finite groups of order at most $n$ by decreasing order. If

\[
|K|>|H|,
\]

then

\[
i_K(H)=0.
\]

If

\[
|K|=|H|,
\]

then

\[
i_K(H)>0
\]

if and only if

\[
K\cong H,
\]

and in that case

\[
i_K(K)=|\operatorname{Aut}(K)|.
\]

Thus the resulting linear system is triangular and determines all multiplicities

\[
\mu_{r,G}(H).
\]

\end{proof}

This shows that CHT strictly refines the centralizer-size moment hierarchy.

\section{Comparison of the observer families}

The observers studied in this paper should not be understood as a single totally ordered chain. Rather, they are several natural families of quadratic observers that retain different kinds of information.

For class-two exponent-$p$ groups $G_\beta$, the higher commuting probabilities

\[
P_m(G_\beta)
\]

are determined by the numbers of totally isotropic subspaces of $V$. In contrast, the star probabilities

\[
\Sigma_{r,t}(G_\beta)
\]

are moments of the rank distribution

\[
N_r^\beta(a)
=
|\{\mathbf v\in V^r:\operatorname{rank}A_{\mathbf v}=a\}|.
\]

Thus $P_m$ and $N_r$ measure different geometric aspects of the alternating map $\beta$. Theorem 6.1 shows that the $P_m$-hierarchy does not determine $N_1$, even on directly indecomposable groups, uniformly for all odd primes $p$.

There are, however, genuine refinement relations among some of the later profiles. The aggregated flag profile of level $r$ refines $N_r$, since it records the ranks of all coordinate subtuples, not only the rank of the full $r$-tuple. The canonical support-image profile $\operatorname{CSI}_r$ further refines the aggregated flag profile: it records the $\operatorname{GL}(W)$-orbit of the actual image subspace

\[
I_\beta(\mathbf v)=\operatorname{im}A_{\mathbf v}\le W^r,
\]

rather than only the dimensions of its coordinate projections.

The Centralizer Hom Transform is of a different nature. It refines the centralizer-size moments by replacing

\[
|C_G(\mathbf x)|^t
\]

with

\[
|\operatorname{Hom}(K,C_G(\mathbf x))|.
\]

For finite test groups $K$, the full transform recovers the isomorphism types of common centralizers with multiplicity. We do not claim that CHT is comparable, in general, with CSI. Rather, it records internal group structure of centralizers, while CSI records codomain-localized linear data in class-two exponent-$p$ groups.

The resulting picture is therefore a ladder of increasingly structured tests in several directions:

\[
\begin{array}{c}
\text{isotropic-subspace data }(P_m),\\[1mm]
\text{rank-support data }(N_r,\Sigma_{r,t}),\\[1mm]
\text{incidence data }(\mathsf{Flag}_r),\\[1mm]
\text{codomain-localized support-image data }(\operatorname{CSI}_r),\\[1mm]
\text{centralizer-type data }(\operatorname{CHT}).
\end{array}
\]

The main point of this paper is that the first two layers are already genuinely distinct, and that the rank-support side admits both a rank-metric-code interpretation and effective codomain-localized refinements.

The results of the paper establish three conclusions.

First, higher commuting probabilities and star profiles encode genuinely different information in class-two exponent-$p$ groups. The former are governed by isotropic-subspace counts, while the latter are governed by rank-support data.

Second, the star rank profiles are not ad hoc invariants: they are precisely the higher rank-support enumerators of the contraction code

\[
C_\beta=\{A_v:v\in V\}.
\]

Third, codomain-localized support-image profiles provide a compact refinement of rank and flag data, computable in fixed-parameter time with respect to the derived rank $\dim W$ and the tuple level $r$.


\section{Further directions}

\subsection{Computational benchmarks}

The $5$-generator class-two exponent-$p$ databases studied by O'Brien and Stanojkovski provide natural testing grounds for $N_r$, flag profiles, and CSI. A separate computational paper will treat these benchmarks with reproducible scripts and certificates.

\subsection{Uniform classification in small derived rank}

The rank-support and CSI profiles suggest a possible classification program for alternating maps

\[
\beta:V\times V\to W
\]

with small $\dim W$. The main problem is to understand when codomain-localized profiles determine $\beta$ up to equivalence.

\subsection{CSI and canonization}

CSI is an invariant, not a canonization procedure in general. However, if the CSI structure on $W^r$ has small stabilizer in $\operatorname{PGL}(W)$, it can be used as a pre-canonization tool. This leads to the study of residual stabilizers of support-image profiles.

\subsection{Centralizer structure beyond class two}

The CHT definition applies to all finite groups. It may be useful for studying groups with identical commuting graphs or identical centralizer-size distributions but different centralizer isomorphism types.

\begin{thebibliography}{99}
\bibitem{gustafson1973} W. H. Gustafson, What is the probability that two group elements commute?, \emph{American Mathematical Monthly} 80 (1973), 1031--1034.
\bibitem{eberhard2015} S. Eberhard, Commuting probabilities of finite groups, \emph{Bulletin of the London Mathematical Society} 47 (2015), 796--808.
\bibitem{levit-shwartz2026} V. E. Levit and R. Shwartz, Higher Commutativity in Finite Groups: Exact Asymptotics and Finite Spectrum, arXiv:2605.02071.
\bibitem{gorla-jurrius-lopez-ravagnani2020} E. Gorla, R. Jurrius, H. H. Lopez, and A. Ravagnani, Rank-Metric Codes and $q$-Polymatroids, \emph{Journal of Algebraic Combinatorics} 52 (2020), 1--19.
\bibitem{sun2023} X. Sun, Faster Isomorphism for $p$-Groups of Class 2 and Exponent $p$, arXiv:2303.15412.
\bibitem{obrien-stanojkovski2024} E. A. O'Brien and M. Stanojkovski, Geometric invariants for $p$-groups of class 2 and exponent $p$, arXiv:2411.19555.
\bibitem{lescot1995} P. Lescot, Isoclinism classes and commutativity degrees of finite groups, \emph{Journal of Algebra} 177 (1995), 847--869.
\bibitem{lovasz1967} L. Lovasz, Operations with structures, \emph{Acta Mathematica Academiae Scientiarum Hungaricae} 18 (1967), 321--328.
\bibitem{atserias-kolaitis-wu2021} A. Atserias, P. G. Kolaitis, and W.-L. Wu, On the Expressive Power of Homomorphism Counts, LICS 2021; arXiv:2101.12733.
\bibitem{wilson2019} J. B. Wilson, The Threshold for Subgroup Profiles to Agree is Logarithmic, \emph{Theory of Computing} 15 (2019), Article 19, 1--25.
\end{thebibliography}
\end{document}
